%! Unidad 4 — El enlace químico: ejercicios resueltos y ejercicios finales
%! Fuente: Fuentes/Quimica/Unidad 4 («Ejercicios resueltos», «Practica y calcula», «Piensa y resuelve»).
%! Enunciados adaptados con sustancias y datos cambiados; soluciones rehechas. Dificultad: \(\star\) baja, \(\star\star\) media, \(\star\star\star\) alta.
%! Electronegatividades de Pauling de la tabla tab:ud3-en (salvo donde el enunciado da otras).

\section{Ejercicios resueltos}\label{sec:ud4-resueltos}

Salvo que se indique otra cosa, usa las electronegatividades de la tabla~\ref{tab:ud3-en} y el criterio de la
tabla~\ref{tab:ud4-den}.

\begin{example}
    \textbf{Lewis, geometría y polaridad.} Determina la estructura de Lewis y la geometría del \ch{PCl3}, el \ch{CS2} y el
    \ch{CH2Cl2}. Usa la configuración electrónica de cada elemento para hallar su valencia covalente. ¿Cuáles son polares?

    \textbf{Solución.} \textbf{1. Datos}: P: \conf{[Ne]\,3s^2\,3p^3} (5 electrones de valencia, covalencia 3); Cl:
    \conf{[Ne]\,3s^2\,3p^5} (7, covalencia 1); C: \conf{[He]\,2s^2\,2p^2} (4, covalencia 4); S: \conf{[Ne]\,3s^2\,3p^4} (6,
    covalencia 2); H: \conf{1s^1} (1, covalencia 1).

    \textbf{2. \ch{PCl3}}: \(D = 5 + 3 \cdot 7 = 26\); \(N = 4 \cdot 8 = 32\); \(C = 6\): \textbf{3 enlaces}. Quedan
    \(26 - 6 = 20\) electrones: tres pares en cada Cl y \textbf{un par solitario en el P}. Con tres pares enlazantes y uno
    solitario, la molécula es una \textbf{pirámide trigonal}; los momentos de los enlaces P--Cl (\(\Delta\mathrm{EN} = 1{,}0\)) no
    se anulan: es \textbf{polar}.
    \[ \chemfig{\charge{90=\:,180=\:,270=\:}{Cl}-\charge{90=\:}{P}(-[6]\charge{0=\:,180=\:,270=\:}{Cl})-\charge{0=\:,90=\:,270=\:}{Cl}} \]

    \textbf{3. \ch{CS2}}: \(D = 4 + 2 \cdot 6 = 16\); \(N = 24\); \(C = 8\): \textbf{4 enlaces}, dos dobles. Quedan 8 electrones,
    dos pares en cada S. Sin pares solitarios en el C, la molécula es \textbf{lineal} y \textbf{apolar} (además,
    \(\Delta\mathrm{EN} = 0\)).
    \[ \chemfig{\charge{90=\:,270=\:}{S}=C=\charge{90=\:,270=\:}{S}} \]

    \textbf{4. \ch{CH2Cl2}}: \(D = 4 + 2 + 14 = 20\); \(N = 8 + 2 \cdot 2 + 2 \cdot 8 = 28\); \(C = 8\): \textbf{4 enlaces}
    simples. Quedan 12 electrones, tres pares en cada Cl. La molécula es \textbf{tetraédrica}, pero los enlaces C--H y C--Cl son
    distintos y sus momentos no se anulan: es \textbf{polar}.
    \[ \chemfig{C(-[2]H)(-[6]H)(-[4]\charge{90=\:,180=\:,270=\:}{Cl})-\charge{0=\:,90=\:,270=\:}{Cl}} \]
\end{example}

\begin{example}
    \textbf{Energía reticular.} Calcula la energía reticular del bromuro de sodio sabiendo que en su formación a partir de sus
    elementos se desprenden \qty{361}{\kilo\joule\per\mol}, y que (en \unit{\kilo\joule\per\mol}): sublimación del Na, \(+108\);
    vaporización y disociación de medio mol de \ch{Br2(l)}, \(+112\); energía de ionización del Na, \(+496\); afinidad electrónica
    del Br, \(-325\). ¿Es mayor o menor que la del \ch{NaCl} (\(\qty{-787}{\kilo\joule\per\mol}\))? ¿Por qué?

    \textbf{Solución.} \textbf{1. Planteamiento}: según el ciclo de Born-Haber, la energía del proceso global es la suma de las de
    todas las etapas:
    \[ E_{\text{formación}} = E_{\text{sub}} + \tfrac{1}{2} E_{\text{dis}} + EI + AE + U \]
    \textbf{2. Cálculo}:
    \[ -361 = 108 + 112 + 496 - 325 + U \;\Rightarrow\; U = -361 - 391 = \qty{-752}{\kilo\joule\per\mol} \]
    \textbf{3. Comparación}: es \textbf{menor} en valor absoluto que la del \ch{NaCl}. Los dos tienen el mismo catión y aniones de
    la misma carga, pero el \ch{Br-} es más grande que el \ch{Cl-}: los iones quedan más lejos y se atraen menos.
\end{example}

\begin{example}
    \textbf{Identificar sustancias por sus propiedades.} Cuatro sólidos, A, B, C y D, tienen estas propiedades:
    \begin{center}
        \begin{tabular}{@{}lcccc@{}}
            \toprule
            & \textbf{Fusión (\unit{\degreeCelsius})} & \textbf{Solubilidad en agua} & \textbf{Conduce en sólido} &
            \textbf{Conduce disuelto o fundido} \\
            \midrule
            A & 1083 & No & Sí & Sí (fundido) \\[3pt]
            B & 113 & No (sí en \ch{CS2}) & No & No \\[3pt]
            C & 1713 & No & No & No \\[3pt]
            D & 714 & Sí & No & Sí \\
            \bottomrule
        \end{tabular}
    \end{center}
    Indica qué tipo de sustancia es cada una.

    \textbf{Solución.} \textbf{A} conduce en estado sólido: es un \textbf{metal} (son los datos del cobre). \textbf{B} funde a baja
    temperatura, no conduce y se disuelve en un disolvente apolar: es una \textbf{sustancia molecular apolar} (el azufre,
    \ch{S8}). \textbf{C} tiene un punto de fusión muy alto, no conduce de ninguna forma y es insoluble: es un \textbf{cristal
    covalente} (el cuarzo, \ch{SiO2}). \textbf{D} funde a temperatura alta, se disuelve en agua y solo conduce cuando sus iones
    quedan libres: es un \textbf{compuesto iónico} (el \ch{MgCl2}).
\end{example}

\begin{example}
    \textbf{Ordenar por punto de ebullición.} Ordena de menor a mayor punto de ebullición estas sustancias: \ch{H2O},
    \ch{CH4}, \ch{C3H8} (propano) y \ch{CH3OH} (metanol).

    \textbf{Solución.} Las cuatro son sustancias moleculares; su punto de ebullición depende de las fuerzas intermoleculares.
    \begin{itemize}
        \item \ch{CH4} y \ch{C3H8} son apolares: solo tienen fuerzas de dispersión, mayores en el propano, cuya molécula es más
        grande.
        \item \ch{CH3OH} y \ch{H2O} tienen grupos O--H y forman enlaces de hidrógeno, mucho más intensos. El agua forma más
        enlaces de hidrógeno por molécula (dos H y dos pares solitarios) que el metanol (un solo H unido a O).
    \end{itemize}
    El orden es \(\ch{CH4} < \ch{C3H8} < \ch{CH3OH} < \ch{H2O}\). Los valores reales son \(-161\), \(-42\), 65 y
    \qty{100}{\degreeCelsius}.
\end{example}

\begin{example}
    \textbf{Del número atómico a las propiedades.} Los elementos X (\(Z = 6\)) e Y (\(Z = 17\)) se combinan entre sí.
    a) ¿Qué compuesto forman? Dibuja su estructura de Lewis. b) ¿Cuál es su geometría? ¿Es polar? c) ¿Qué fuerzas hay entre sus
    moléculas? ¿Será soluble en agua?

    \textbf{Solución.} a) X: \conf{1s^2\,2s^2\,2p^2} es el carbono (covalencia 4) e Y: \conf{1s^2\,2s^2\,2p^6\,3s^2\,3p^5}, el cloro
    (covalencia 1). Los dos son no metales: forman un compuesto covalente, el \ch{CCl4}:
    \[ \chemfig{C(-[2]\charge{0=\:,90=\:,180=\:}{Cl})(-[4]\charge{90=\:,180=\:,270=\:}{Cl})(-[6]\charge{0=\:,180=\:,270=\:}{Cl})-\charge{0=\:,90=\:,270=\:}{Cl}} \]
    b) El C tiene cuatro pares enlazantes y ninguno solitario: geometría \textbf{tetraédrica}. Los enlaces C--Cl son polares
    (\(\Delta\mathrm{EN} = \num{0.6}\)), pero los cuatro momentos, iguales y dirigidos hacia los vértices de un tetraedro, se
    anulan: la molécula es \textbf{apolar}. c) Entre sus moléculas solo hay \textbf{fuerzas de dispersión}, bastante intensas por su
    tamaño (es un líquido que hierve a \qty{77}{\degreeCelsius}). Al ser apolar, \textbf{no es soluble en agua}; de hecho, es un
    disolvente de sustancias apolares.
\end{example}

%========================================================================================
\section{Ejercicios de aplicación}\label{sec:ud4-aplicacion}
%========================================================================================

\subsection{Enlace iónico}

\begin{exercise}
    \(\star\star\) Predice el ion más probable que formarán el oxígeno, el litio, el estroncio y el yodo. Usa la regla del octeto
    y sus configuraciones electrónicas.
    \begin{solu}
        O: \conf{[He]\,2s^2\,2p^4}, gana 2: \textcolor{red}{\ch{O^2-}}. Li: \conf{[He]\,2s^1}, pierde 1: \textcolor{red}{\ch{Li+}}.
        Sr: \conf{[Kr]\,5s^2}, pierde 2: \textcolor{red}{\ch{Sr^2+}}. I: \conf{[Kr]\,4d^{10}\,5s^2\,5p^5}, gana 1:
        \textcolor{red}{\ch{I-}}. Todos quedan con la configuración de un gas noble.
    \end{solu}
\end{exercise}

\begin{exercise}
    \(\star\) De los compuestos \ch{KCl}, \ch{ClF3} y \ch{HI}, ¿cuál tiene enlace iónico? Justifica la respuesta.
    \begin{solu}
        \textcolor{red}{El \ch{KCl}}: está formado por un metal y un no metal (\(\Delta\mathrm{EN} = 3{,}2 - 0{,}8 = 2{,}4\)). El
        \ch{ClF3} y el \ch{HI} están formados solo por no metales y son covalentes.
    \end{solu}
\end{exercise}

\begin{exercise}
    \(\star\star\) De estas afirmaciones sobre el enlace iónico, solo una es falsa. ¿Cuál? a) Se da entre un metal y un no metal.
    b) Los compuestos iónicos conducen la electricidad en estado sólido. c) Se forma entre elementos con mucha diferencia de
    electronegatividad. d) Muchos compuestos iónicos son solubles en disolventes polares.
    \begin{solu}
        \textcolor{red}{La b}: en estado sólido los iones ocupan posiciones fijas en la red y no pueden moverse. Solo conducen
        fundidos o disueltos.
    \end{solu}
\end{exercise}

\begin{exercise}
    \(\star\star\) De estas afirmaciones sobre los sólidos iónicos, solo una es verdadera. ¿Cuál? a) Tienen puntos de fusión bajos.
    b) Son más solubles en disolventes apolares que en agua. c) Conducen la electricidad fundidos. d) Son dúctiles y maleables.
    \begin{solu}
        \textcolor{red}{La c}. Tienen puntos de fusión altos, se disuelven en disolventes polares y son frágiles.
    \end{solu}
\end{exercise}

\begin{exercise}
    \(\star\) Si combinamos un elemento X con tres electrones en su capa de valencia con un elemento Y con seis, ¿cuál será la
    fórmula del compuesto iónico que forman?
    \begin{solu}
        X pierde 3 electrones (\ch{X^3+}) e Y gana 2 (\ch{Y^2-}). Para que el compuesto sea neutro, 2 \ch{X^3+} (6 cargas
        positivas) por cada 3 \ch{Y^2-} (6 negativas): \textcolor{red}{\ch{X2Y3}}, como el \ch{Al2O3}.
    \end{solu}
\end{exercise}

\begin{exercise}
    \(\star\star\) Investiga y explica a tus compañeros qué son la energía reticular y el índice de coordinación de un cristal
    iónico. ¿De qué depende cada uno?
    \begin{solu}
        \textcolor{red}{Energía reticular: la que se desprende al formarse un mol de cristal a partir de sus iones gaseosos}; aumenta
        con la carga de los iones y al disminuir su tamaño. \textcolor{red}{Índice de coordinación: el número de iones de signo
        contrario que rodean a cada ion}; depende de la relación entre los tamaños del catión y el anión y de la proporción entre
        ellos (6 en el \ch{NaCl}, 8 en el \ch{CsCl}, cuyo catión es mayor).
    \end{solu}
\end{exercise}

\begin{exercise}
    \(\star\star\) Explica cómo se forma el enlace iónico en: a) \ch{Li2O}; b) \ch{BaF2}; c) \ch{AlF3}.
    \begin{solu}
        a) Cada Li pierde un electrón y el O gana dos: \textcolor{red}{dos \ch{Li+} por cada \ch{O^2-}}. b) El Ba pierde dos
        electrones y cada F gana uno: \textcolor{red}{un \ch{Ba^2+} por cada dos \ch{F-}}. c) El Al pierde tres electrones:
        \textcolor{red}{un \ch{Al^3+} por cada tres \ch{F-}}. Todos los iones quedan con configuración de gas noble y se atraen
        formando una red cristalina.
    \end{solu}
\end{exercise}

\begin{exercise}
    \(\star\star\) Explica por qué los elementos de los grupos 1 y 2, al unirse con los de los grupos 16 y 17, forman sin
    dificultad enlaces iónicos. Pon ejemplos de las uniones posibles.
    \begin{solu}
        \textcolor{red}{Los de los grupos 1 y 2 tienen energías de ionización bajas} y pierden fácilmente 1 o 2 electrones;
        \textcolor{red}{los de los grupos 16 y 17 son muy electronegativos} y ganan 2 o 1. La diferencia de electronegatividad es
        muy grande. Ejemplos: \ch{NaCl}, \ch{KBr}, \ch{Na2O}, \ch{K2S}, \ch{MgF2}, \ch{CaCl2}, \ch{MgO}, \ch{BaS}.
    \end{solu}
\end{exercise}

\begin{exercise}
    \(\star\) Una de las características de los cristales iónicos es que son muy frágiles. ¿A qué se debe? Justifica la respuesta
    con un ejemplo.
    \begin{solu}
        \textcolor{red}{Al golpear el cristal, una capa de iones se desplaza respecto de otra y quedan enfrentados iones del mismo
        signo, que se repelen}: el cristal se parte. Por ejemplo, un cristal de sal gema se rompe en fragmentos de caras planas al
        golpearlo con un martillo (figura~\ref{fig:ud4-deformacion}).
    \end{solu}
\end{exercise}

\begin{exercise}
    \(\star\star\) ¿Qué compuestos formarán estas parejas de elementos? a) Ca y Cl; b) Na y O; c) Rb e I; d) Mg y S. Representa
    la formación de los iones con diagramas de Lewis.
    \begin{solu}
        a) \textcolor{red}{\ch{CaCl2}}: \(\ch{Ca^2+}\) y dos \(\ch{Cl-}\). b) \textcolor{red}{\ch{Na2O}}: dos \ch{Na+} y un
        \ch{O^2-}. c) \textcolor{red}{\ch{RbI}}: \ch{Rb+} y \ch{I-}. d) \textcolor{red}{\ch{MgS}}: \ch{Mg^2+} y \ch{S^2-}. En los
        diagramas, el metal pierde sus puntos y el no metal queda con ocho, entre corchetes y con su carga; por ejemplo, para el
        \ch{MgS}: \(\ch{Mg^2+}\;[\chemfig{\charge{0=\:,90=\:,180=\:,270=\:}{S}}]^{2-}\).
    \end{solu}
\end{exercise}

\begin{exercise}
    \(\star\star\) Calcula la energía que se desprende al formarse un mol de cloruro de cesio a partir de sus elementos, con estos
    datos (en \unit{\kilo\joule\per\mol}): sublimación del Cs, \(+76\); disociación de medio mol de \ch{Cl2}, \(+121\); ionización
    del Cs, \(+376\); afinidad electrónica del Cl, \(-349\); energía reticular del \ch{CsCl}, \(-659\).
    \begin{solu}
        \(E = 76 + 121 + 376 - 349 - 659 = \textcolor{red}{\qty{-435}{\kilo\joule\per\mol}}\).
    \end{solu}
\end{exercise}

\begin{exercise}
    \(\star\star\) Ordena de mayor a menor energía reticular (en valor absoluto): \ch{KF}, \ch{CaO}, \ch{KI} y \ch{LiF}.
    \begin{solu}
        El \ch{CaO} tiene iones de carga \(\pm 2\): es el mayor. Entre los demás, con cargas \(\pm 1\), cuanto más pequeños los
        iones, mayor la energía: \textcolor{red}{\(\ch{CaO} > \ch{LiF} > \ch{KF} > \ch{KI}\)} (\num{3414}, \num{1037}, 821 y
        \qty{649}{\kilo\joule\per\mol}).
    \end{solu}
\end{exercise}

\subsection{Enlace metálico}

\begin{exercise}
    \(\star\star\) Teniendo en cuenta las configuraciones electrónicas del estroncio y del indio, determina cuántos electrones
    cede cada uno a la nube electrónica.
    \begin{solu}
        Sr: \conf{[Kr]\,5s^2} \(\rightarrow\) \textcolor{red}{cede 2} y queda como \ch{Sr^2+}. In: \conf{[Kr]\,4d^{10}\,5s^2\,5p^1}
        \(\rightarrow\) \textcolor{red}{cede 3} y queda como \ch{In^3+} (\conf{[Kr]\,4d^{10}}, con el subnivel \(d\) lleno).
    \end{solu}
\end{exercise}

\begin{exercise}
    \(\star\star\) El enlace metálico une fuertemente a los átomos. ¿Cómo se explica entonces que los metales sean dúctiles y
    maleables sin romperse?
    \begin{solu}
        \textcolor{red}{El enlace metálico no tiene direcciones fijas}: los cationes están unidos a la nube de electrones en su
        conjunto, no a unos vecinos concretos. Cuando una capa se desliza sobre otra, los electrones se redistribuyen y siguen
        uniendo a los cationes, sin que aparezcan repulsiones entre iones del mismo signo.
    \end{solu}
\end{exercise}

\begin{exercise}
    \(\star\) ¿Por qué los cables eléctricos son de cobre, pero van recubiertos de plástico?
    \begin{solu}
        \textcolor{red}{El cobre conduce muy bien la electricidad, gracias a sus electrones deslocalizados; el plástico es una
        sustancia covalente cuyos electrones están localizados en los enlaces}, así que es aislante y evita descargas y
        cortocircuitos.
    \end{solu}
\end{exercise}

\begin{exercise}
    \(\star\) A excepción del mercurio, todos los metales son sólidos a temperatura ambiente. Investiga por qué el mercurio es
    líquido.
    \begin{solu}
        \textcolor{red}{Respuesta abierta.} \textcolor{blue}{El mercurio tiene completos los subniveles \conf{5d^{10}} y
        \conf{6s^2}, y sus electrones 6s, por ser un átomo tan pesado, están muy atraídos por el núcleo (un efecto que se explica con
        la teoría de la relatividad): los cede con dificultad a la nube electrónica, el enlace metálico es muy débil y funde a
        \qty{-39}{\degreeCelsius}.}
    \end{solu}
\end{exercise}

\begin{exercise}
    \(\star\) ¿Por qué los metales son buenos conductores de la electricidad y del calor?
    \begin{solu}
        \textcolor{red}{Porque tienen electrones deslocalizados que se mueven libremente por toda la red}: al aplicar una diferencia
        de potencial se desplazan de forma ordenada (corriente eléctrica), y al calentar una zona transmiten rápidamente la energía
        al resto.
    \end{solu}
\end{exercise}

\begin{exercise}
    \(\star\star\) Para explicar el enlace metálico hemos usado el modelo de la nube electrónica. Buscad información sobre la
    teoría de bandas y preparad un breve trabajo que explique por qué unos materiales son conductores, otros semiconductores y
    otros aislantes.
    \begin{solu}
        \textcolor{red}{Respuesta abierta.} Debe incluir que en un sólido los orbitales de todos los átomos forman
        \textcolor{red}{bandas de energía}; la conducción exige que los electrones puedan pasar a una banda con niveles vacíos
        (\textcolor{red}{banda de conducción}). En los metales la banda de valencia está semillena o solapada con la de
        conducción; en los semiconductores las separa un salto pequeño, y en los aislantes, uno grande
        (figura~\ref{fig:ud4-bandas}).
    \end{solu}
\end{exercise}

\begin{exercise}
    \(\star\) Investiga: ¿por qué los metales alcalinos no cumplen las propiedades típicas de los metales en cuanto a densidad,
    dureza y brillo?
    \begin{solu}
        \textcolor{red}{Cada átomo aporta un solo electrón a la nube y sus átomos son muy grandes}: el enlace metálico es débil,
        así que son blandos (se cortan con un cuchillo) y funden a temperaturas bajas, y su densidad es pequeña (el litio, el sodio
        y el potasio flotan en el agua). \textcolor{red}{Su brillo dura muy poco porque reaccionan enseguida con el oxígeno y la
        humedad del aire}; por eso se guardan en aceite.
    \end{solu}
\end{exercise}

\begin{exercise}
    \(\star\star\) De estas afirmaciones sobre los metales, solo una es falsa. ¿Cuál? a) Son buenos conductores del calor y la
    electricidad. b) Son frágiles. c) Forman aleaciones. d) Tienen brillo característico.
    \begin{solu}
        \textcolor{red}{La b}: son tenaces, dúctiles y maleables, porque las capas de cationes se pueden deslizar sin que el metal
        se rompa.
    \end{solu}
\end{exercise}

\subsection{Enlace covalente}

\begin{exercise}
    \(\star\star\) Dibuja la estructura de Lewis del \ch{C2H4}, el \ch{SiH4}, el \ch{CH3OH}, el \ch{PBr3} y el \ch{H2S}.
    Justifica en cada caso la valencia covalente de los átomos con su configuración electrónica.
    \begin{solu}
        C y Si (\conf{ns^2\,np^2}): covalencia 4. O y S (\conf{ns^2\,np^4}): 2. P (\conf{3s^2\,3p^3}): 3. Br
        (\conf{4s^2\,4p^5}) y H (\conf{1s^1}): 1.
        \begin{center}
            \textcolor{red}{\chemfig{C(-[:150]H)(-[:210]H)=C(-[:30]H)-[:-30]H}} \quad
            \textcolor{red}{\chemfig{Si(-[2]H)(-[4]H)(-[6]H)-H}} \quad
            \textcolor{red}{\chemfig{C(-[2]H)(-[4]H)(-[6]H)-\charge{90=\:,270=\:}{O}-H}} \\[6pt]
            \textcolor{red}{\chemfig{\charge{90=\:,180=\:,270=\:}{Br}-\charge{90=\:}{P}(-[6]\charge{0=\:,180=\:,270=\:}{Br})-\charge{0=\:,90=\:,270=\:}{Br}}} \quad
            \textcolor{red}{\chemfig{H-\charge{90=\:,270=\:}{S}-H}}
        \end{center}
        En el \ch{C2H4} los carbonos se unen con un enlace doble.
    \end{solu}
\end{exercise}

\begin{exercise}
    \(\star\star\) ¿Cuál es la valencia covalente de los átomos en las moléculas \ch{HF}, \ch{O2} y \ch{N2}? Ayúdate de sus
    estructuras de Lewis.
    \begin{solu}
        \textcolor{red}{H y F: 1} (enlace simple, \chemfig{H-\charge{0=\:,90=\:,270=\:}{F}}). \textcolor{red}{O: 2} (enlace doble,
        \chemfig{\charge{90=\:,270=\:}{O}=\charge{90=\:,270=\:}{O}}). \textcolor{red}{N: 3} (enlace triple,
        \chemfig{\charge{180=\:}{N}~\charge{0=\:}{N}}).
    \end{solu}
\end{exercise}

\begin{exercise}
    \(\star\star\) Explica las diferencias entre el enlace covalente, el iónico y el metálico: qué elementos se unen, qué hacen
    con sus electrones y qué tipo de estructura forman.
    \begin{solu}
        \textcolor{red}{Iónico}: metal y no metal; el metal cede electrones al no metal; red de iones. \textcolor{red}{Covalente}:
        no metales; comparten pares de electrones, localizados entre los átomos; moléculas o redes de átomos.
        \textcolor{red}{Metálico}: átomos de metales; ceden sus electrones de valencia a una nube común, deslocalizada; red de
        cationes inmersa en los electrones.
    \end{solu}
\end{exercise}

\begin{exercise}
    \(\star\star\) El carbono y el oxígeno forman dos óxidos, el \ch{CO} y el \ch{CO2}. Dibuja sus estructuras de Lewis. ¿Qué
    tiene de particular uno de los enlaces del \ch{CO}?
    \begin{solu}
        \ch{CO2}: \textcolor{red}{\chemfig{\charge{90=\:,270=\:}{O}=C=\charge{90=\:,270=\:}{O}}}.
        \ch{CO}: \(D = 10\), \(N = 16\), \(C = 6\): \textcolor{red}{3 enlaces, un triple}, con un par solitario en cada átomo:
        \textcolor{red}{\chemfig{\charge{180=\:}{C}~\charge{0=\:}{O}}}. Como el C solo tiene 4 electrones de valencia,
        \textcolor{red}{uno de los tres pares lo aporta solo el oxígeno: es un enlace dativo}.
    \end{solu}
\end{exercise}

\begin{exercise}
    \(\star\star\) De estas afirmaciones sobre el enlace covalente, solo una es falsa. ¿Cuál? a) Puede dar lugar a moléculas o a
    redes de átomos. b) Los átomos se ceden electrones. c) Puede formar sustancias simples. d) Los electrones del enlace están
    localizados entre los átomos que se unen.
    \begin{solu}
        \textcolor{red}{La b}: en el enlace covalente los átomos comparten electrones; cederlos es propio del enlace iónico (y del
        metálico).
    \end{solu}
\end{exercise}

\begin{exercise}
    \(\star\star\) Los elementos X (\(Z = 14\)) e Y (\(Z = 9\)) se combinan. ¿Qué compuesto forman? Dibuja su estructura de Lewis.
    \begin{solu}
        X: \conf{1s^2\,2s^2\,2p^6\,3s^2\,3p^2} (silicio, covalencia 4); Y: \conf{1s^2\,2s^2\,2p^5} (flúor, covalencia 1). Forman
        \textcolor{red}{\ch{SiF4}}, con el Si en el centro unido a cuatro F por enlaces simples y tres pares solitarios en cada F:
        \textcolor{red}{\chemfig{Si(-[2]\charge{0=\:,90=\:,180=\:}{F})(-[4]\charge{90=\:,180=\:,270=\:}{F})(-[6]\charge{0=\:,180=\:,270=\:}{F})-\charge{0=\:,90=\:,270=\:}{F}}}.
    \end{solu}
\end{exercise}

\begin{exercise}
    \(\star\star\) Un compañero ha dibujado así varias estructuras de Lewis: a) \ch{Cl2}, con un enlace doble y dos pares
    solitarios en cada Cl; b) \ch{O2}, con un enlace simple y tres pares solitarios en cada O; c) \ch{HBr}, con un enlace simple,
    un par solitario en el H y tres en el Br; d) \ch{CH4}, con cuatro enlaces simples y sin pares solitarios. ¿Cuáles son
    incorrectas? Corrígelas.
    \begin{solu}
        \textcolor{red}{Solo la d es correcta.} a) Cada Cl tendría 8 electrones, pero se han usado 12 en total y el Cl tiene
        covalencia 1: \chemfig{\charge{90=\:,180=\:,270=\:}{Cl}-\charge{0=\:,90=\:,270=\:}{Cl}}. b) Así cada O solo tendría 7
        electrones: hace falta un enlace doble, \chemfig{\charge{90=\:,270=\:}{O}=\charge{90=\:,270=\:}{O}}. c) El H solo puede
        tener 2 electrones (el del enlace): \chemfig{H-\charge{0=\:,90=\:,270=\:}{Br}}.
    \end{solu}
\end{exercise}

\begin{exercise}
    \(\star\star\) El hidrógeno forma compuestos con casi todos los elementos. Si se une con O, Cl, Br, C y Si, ¿cuál de los
    compuestos tendrá enlaces más polares? ¿Y menos polares?
    \begin{solu}
        \(\Delta\mathrm{EN}\): O--H 1,2; Cl--H 1,0; Br--H 0,8; C--H 0,4; Si--H 0,3. \textcolor{red}{Más polares: los O--H del
        \ch{H2O}}. \textcolor{red}{Menos polares: los Si--H del \ch{SiH4}}, en los que además el polo negativo está en el H.
    \end{solu}
\end{exercise}

\begin{exercise}
    \(\star\) Ordena de menor a mayor polaridad de enlace: \ch{O2}, \ch{NH3} y \ch{HF}.
    \begin{solu}
        \(\Delta\mathrm{EN}\): 0; 0,8; 1,8: \textcolor{red}{\(\ch{O2} < \ch{NH3} < \ch{HF}\)}.
    \end{solu}
\end{exercise}

\begin{exercise}
    \(\star\star\) ¿Cuál de estas afirmaciones sobre los enlaces polares es verdadera? a) Los electrones se reparten por igual.
    b) Los dos átomos son iguales. c) Los dos átomos tienen distinta electronegatividad. d) Todas son verdaderas.
    \begin{solu}
        \textcolor{red}{La c}: el átomo más electronegativo atrae más los electrones compartidos, que no se reparten por igual.
    \end{solu}
\end{exercise}

\begin{exercise}
    \(\star\star\) Ordena de menor a mayor polaridad los enlaces H--Na, H--O, H--Cl y H--P, con estas electronegatividades:
    H 2,2; Na 0,9; O 3,4; Cl 3,2; P 2,2. ¿En cuál de ellos tiene el hidrógeno carga parcial negativa?
    \begin{solu}
        \(\Delta\mathrm{EN}\): H--P 0; H--Cl 1,0; H--O 1,2; H--Na 1,3. \textcolor{red}{\(\text{H--P} < \text{H--Cl} < \text{H--O} <
        \text{H--Na}\)}. \textcolor{red}{En el H--Na}: el H es más electronegativo que el sodio, y en el hidruro de sodio, \ch{NaH},
        el hidrógeno está como ion \ch{H-}.
    \end{solu}
\end{exercise}

\subsection{Sustancias covalentes moleculares}

\begin{exercise}
    \(\star\star\) Con la configuración electrónica, la valencia covalente y la estructura de Lewis, determina si estas sustancias
    tienen geometría tetraédrica, de pirámide trigonal o triangular plana: \ch{CF4}, \ch{PCl3}, \ch{BBr3} y \ch{H3O+}.
    \begin{solu}
        \ch{CF4}: 4 pares enlazantes, \textcolor{red}{tetraédrica}. \ch{PCl3}: 3 enlazantes y 1 solitario, \textcolor{red}{pirámide
        trigonal}. \ch{BBr3}: 3 enlazantes (octeto incompleto), \textcolor{red}{triangular plana}. \ch{H3O+}: 3 enlazantes y 1
        solitario, \textcolor{red}{pirámide trigonal}.
    \end{solu}
\end{exercise}

\begin{exercise}
    \(\star\star\) Si el enlace Si--F es polar, ¿por qué la molécula de \ch{SiF4} no tiene momento dipolar?
    \begin{solu}
        \textcolor{red}{Porque es tetraédrica}: los cuatro momentos de enlace son iguales y apuntan hacia los vértices de un
        tetraedro regular, así que su suma vectorial es cero.
    \end{solu}
\end{exercise}

\begin{exercise}
    \(\star\star\) Dibuja la estructura de Lewis del \ch{GeCl4}, el \ch{H2S}, el \ch{PH3}, el \ch{N2H4} (hidracina) y el
    \ch{CaCl2}, y predice su geometría. ¿Hay alguno al que no se le pueda aplicar la TRPECV?
    \begin{solu}
        \ch{GeCl4}: \textcolor{red}{tetraédrica}. \ch{H2S}: dos pares enlazantes y dos solitarios, \textcolor{red}{angular}.
        \ch{PH3}: \textcolor{red}{pirámide trigonal}. \ch{N2H4}: \chemfig{H-\charge{90=\:}{N}(-[6]H)-\charge{90=\:}{N}(-[6]H)-H};
        \textcolor{red}{cada N es el vértice de una pirámide trigonal}. \textcolor{red}{El \ch{CaCl2} es iónico}: no forma
        moléculas, sino una red de iones \ch{Ca^2+} y \ch{Cl-}, y la TRPECV no se le aplica.
    \end{solu}
\end{exercise}

\begin{exercise}
    \(\star\star\) Dibuja la estructura de Lewis del \ch{AlCl3} y del \ch{PF3}. ¿Qué diferencia hay entre ellas? ¿Tienen la
    misma geometría?
    \begin{solu}
        \textcolor{red}{En el \ch{AlCl3} el aluminio forma tres enlaces y se queda con 6 electrones (octeto incompleto); en el
        \ch{PF3} el fósforo forma tres enlaces y le queda un par solitario}, así que cumple el octeto.
        \textcolor{red}{No}: el \ch{AlCl3} es triangular plano (en estado gaseoso) y el \ch{PF3}, una pirámide trigonal.
    \end{solu}
\end{exercise}

\begin{exercise}
    \(\star\star\) ¿Cuáles de estas sustancias están formadas por moléculas con momento dipolar? a) \ch{PF3}; b) \ch{CCl4};
    c) \ch{MgCl2}; d) \ch{CH3Cl}; e) \ch{H2O}. Dibuja sus estructuras de Lewis para justificarlo.
    \begin{solu}
        \textcolor{red}{a) \ch{PF3} (piramidal), d) \ch{CH3Cl} (tetraédrica con un átomo distinto) y e) \ch{H2O} (angular)}. El
        \ch{CCl4} es tetraédrico y simétrico: \(\vec{\mu} = 0\). El \ch{MgCl2} es iónico: no forma moléculas.
    \end{solu}
\end{exercise}

\begin{exercise}
    \(\star\star\) ¿Pueden formar enlaces de hidrógeno las sustancias \ch{CH3OH}, \ch{H2S}, \ch{NH3}, \ch{HCl} y \ch{HF}? Razona
    la respuesta.
    \begin{solu}
        \textcolor{red}{Sí el \ch{CH3OH}, el \ch{NH3} y el \ch{HF}}, que tienen H unidos a O, N o F. \textcolor{red}{No el \ch{H2S}
        ni el \ch{HCl}}: el S y el Cl no son lo bastante electronegativos y pequeños (entre sus moléculas hay fuerzas
        dipolo-dipolo).
    \end{solu}
\end{exercise}

\begin{exercise}
    \(\star\) ¿Tienen dipolos permanentes las moléculas de \ch{HBr}, \ch{H2S} y \ch{N2}?
    \begin{solu}
        \textcolor{red}{Sí el \ch{HBr} y el \ch{H2S}}: tienen enlaces polares y no son simétricas (el \ch{H2S} es angular).
        \textcolor{red}{El \ch{N2} no}: su enlace es apolar.
    \end{solu}
\end{exercise}

\begin{exercise}
    \(\star\star\) El ADN está formado por dos cadenas de nucleótidos enrolladas en una doble hélice. Investiga qué fuerzas
    mantienen unidas las dos cadenas y por qué es importante que sean relativamente débiles.
    \begin{solu}
        \textcolor{red}{Enlaces de hidrógeno entre las bases nitrogenadas} (dos entre adenina y timina, tres entre guanina y
        citosina), además de fuerzas de dispersión entre las bases apiladas. \textcolor{red}{Al ser débiles, las dos hebras se
        pueden separar sin romper enlaces covalentes} cuando la célula copia el ADN o lee sus genes, y volver a unirse después.
    \end{solu}
\end{exercise}

\begin{exercise}
    \(\star\star\) El carbono y el silicio están en el mismo grupo. Al unirse con oxígeno forman \ch{CO2}, un gas, y \ch{SiO2},
    un sólido que funde a más de \qty{1700}{\degreeCelsius}. ¿A qué se debe esta diferencia?
    \begin{solu}
        \textcolor{red}{El \ch{CO2} es una sustancia molecular}: el C forma dobles enlaces con dos O y la molécula, apolar, solo se
        une a las demás por fuerzas de dispersión. \textcolor{red}{El \ch{SiO2} es un cristal covalente}: cada Si se une con
        enlaces simples a cuatro O, y cada O a dos Si, formando una red gigante; para fundirlo hay que romper enlaces covalentes.
    \end{solu}
\end{exercise}

\begin{exercise}
    \(\star\star\) Haz un resumen de todas las fuerzas intermoleculares: entre qué moléculas actúan, a qué se deben y cuál es su
    intensidad relativa.
    \begin{solu}
        \textcolor{red}{Dispersión}: entre todas las moléculas; dipolos instantáneos e inducidos; débiles, crecen con el tamaño de
        la molécula. \textcolor{red}{Dipolo-dipolo}: entre moléculas polares; atracción entre dipolos permanentes; algo más
        intensas. \textcolor{red}{Enlace de hidrógeno}: entre moléculas con H unido a F, O o N; las más intensas, aunque unas diez
        veces más débiles que un enlace covalente.
    \end{solu}
\end{exercise}

\begin{exercise}
    \(\star\star\) La sílice (\ch{SiO2}) es el componente principal del cuarzo y del vidrio. Aunque su composición es la misma,
    sus estructuras son muy diferentes. Investiga las estructuras, los usos y las propiedades de estos dos materiales.
    \begin{solu}
        \textcolor{red}{Respuesta abierta.} \textcolor{blue}{El cuarzo es un cristal covalente: tetraedros \ch{SiO4} ordenados que se
        repiten; es muy duro (7 en la escala de Mohs), funde a \qty{1713}{\degreeCelsius} y es piezoeléctrico (se usa en relojes y
        circuitos). El vidrio tiene los mismos tetraedros, pero unidos al azar (estructura amorfa); se reblandece en un intervalo
        de temperaturas, lo que permite moldearlo, y se le añaden óxidos de sodio y calcio para bajar su temperatura de trabajo.}
    \end{solu}
\end{exercise}

\begin{exercise}
    \(\star\star\) Para fabricar móviles, tabletas y ordenadores es imprescindible el coltán, un mineral escaso. Investiga qué
    se obtiene de él y para qué se usa, dónde están sus principales reservas y qué consecuencias tiene su extracción para las
    personas que trabajan en las minas y para el medio ambiente.
    \begin{solu}
        \textcolor{red}{Respuesta abierta.} \textcolor{blue}{Del coltán se extraen el tántalo y el niobio. El tántalo se usa en
        condensadores muy pequeños y estables de los aparatos electrónicos. Gran parte procede de la República Democrática del
        Congo y de Ruanda; su extracción ha financiado a grupos armados, se ha hecho con trabajo infantil y en condiciones muy
        peligrosas, y ha provocado deforestación y daños en los hábitats del gorila. Hoy existen sistemas de certificación de
        «minerales libres de conflicto».}
    \end{solu}
\end{exercise}

%========================================================================================
\section{Ejercicios de ampliación}\label{sec:ud4-ampliacion}
%========================================================================================

\subsection{El enlace químico}

\begin{exercise}
    \(\star\star\star\) Determina qué tipo de enlace tienen el \ch{F2}, el \ch{NH3}, el \ch{MgCl2} y el \ch{Ni}. Dibuja la estructura
    de Lewis de los dos primeros. a) ¿Cuál de ellos es polar? b) ¿Cuál es apolar? c) ¿Cuáles son solubles en agua? d) ¿Conducirá
    alguno la electricidad? ¿En qué condiciones?
    \begin{solu}
        \ch{F2}: \textcolor{red}{covalente} (\chemfig{\charge{90=\:,180=\:,270=\:}{F}-\charge{0=\:,90=\:,270=\:}{F}});
        \ch{NH3}: \textcolor{red}{covalente} (\chemfig{H-\charge{90=\:}{N}(-[6]H)-H}); \ch{MgCl2}: \textcolor{red}{iónico}; \ch{Ni}:
        \textcolor{red}{metálico}. a) \textcolor{red}{\ch{NH3}}. b) \textcolor{red}{\ch{F2}}. c) \textcolor{red}{\ch{NH3}} (forma
        enlaces de hidrógeno con el agua) \textcolor{red}{y \ch{MgCl2}}. d) \textcolor{red}{El níquel, siempre; el \ch{MgCl2},
        fundido o disuelto.}
    \end{solu}
\end{exercise}

\begin{exercise}
    \(\star\star\star\) Representa la estructura de Lewis del \ch{C2H4}, el \ch{HCN}, el \ch{H2O2} y el \ch{N2}, e indica cuáles
    tienen algún enlace doble y cuáles alguno triple.
    \begin{solu}
        \textcolor{red}{Doble: \ch{C2H4}} (C=C). \textcolor{red}{Triple: \ch{HCN}} (C\(\equiv\)N) \textcolor{red}{y \ch{N2}}.
        \textcolor{red}{El \ch{H2O2} solo tiene enlaces simples}:
        \chemfig{H-\charge{90=\:,270=\:}{O}-\charge{90=\:,270=\:}{O}-H}.
    \end{solu}
\end{exercise}

\begin{exercise}
    \(\star\star\star\) En el \ch{H2O}, el \ch{NH4+}, el \ch{BCl3} y el \ch{PCl3}, ¿en cuáles tiene el átomo central algún par de
    electrones no enlazante? ¿Cómo influye en su geometría?
    \begin{solu}
        \textcolor{red}{En el \ch{H2O} (dos pares) y en el \ch{PCl3} (uno)}. En el \ch{NH4+} el par solitario del N se ha usado
        en el enlace dativo con el \ch{H+}, y el B del \ch{BCl3} solo tiene 3 electrones. Los pares solitarios hacen que el agua
        sea angular y el \ch{PCl3}, piramidal; los otros dos son tetraédrico (\ch{NH4+}) y triangular plano (\ch{BCl3}).
    \end{solu}
\end{exercise}

\begin{exercise}
    \(\star\star\star\) ¿Cuál de estos compuestos cumple la regla del octeto? a) \ch{PCl5}; b) \ch{XeF4}; c) \ch{CF4}; d) \ch{BF3}.
    \begin{solu}
        \textcolor{red}{Solo el \ch{CF4}}. El P del \ch{PCl5} tiene 10 electrones y el Xe del \ch{XeF4}, 12 (octetos expandidos); el
        B del \ch{BF3} tiene 6 (octeto incompleto).
    \end{solu}
\end{exercise}

\begin{exercise}
    \(\star\star\star\) ¿Cuál de estas especies no cumple la regla del octeto? a) \ch{CHCl3}; b) \ch{NH4+}; c) \ch{BeH2};
    d) \ch{CO2}.
    \begin{solu}
        \textcolor{red}{El \ch{BeH2}}: el Be forma dos enlaces y queda con 4 electrones.
    \end{solu}
\end{exercise}

\begin{exercise}
    \(\star\star\star\) Cuatro elementos, A, B, C y D, tienen electronegatividades 4,0; 3,2; 2,6 y 1,0, respectivamente. Forman
    los compuestos BA, CA, DA y DC. a) ¿Cuál tendrá mayor carácter covalente? b) ¿Cuál será más polar? c) ¿Serán todos solubles en
    agua? d) ¿Cuál tendrá mayor punto de fusión?
    \begin{solu}
        \(\Delta\mathrm{EN}\): BA 0,8; CA 1,4; DA 3,0; DC 1,6. a) \textcolor{red}{BA}, el de menor diferencia. b)
        \textcolor{red}{DA}, que es prácticamente iónico. c) \textcolor{red}{No tienen por qué}: DA (iónico) y los compuestos
        covalentes polares pueden serlo, pero depende de la geometría de las moléculas y de la energía reticular. d)
        \textcolor{red}{DA}: es el compuesto iónico, con una red de iones fuertemente unidos.
    \end{solu}
\end{exercise}

\begin{exercise}
    \(\star\star\star\) Razona si estas afirmaciones son ciertas: a) El \ch{CH4} es polar y el \ch{NH3}, apolar. b) El \ch{CH4} es
    soluble en agua y el \ch{KCl}, insoluble. c) El \ch{KCl} tiene un punto de fusión mayor que el \ch{CH4}.
    \begin{solu}
        a) \textcolor{red}{Falsa}: al revés; el \ch{CH4} es tetraédrico y simétrico, y el \ch{NH3}, piramidal y polar.
        b) \textcolor{red}{Falsa}: al revés; el \ch{CH4} es apolar y el \ch{KCl}, iónico. c) \textcolor{red}{Verdadera}: hay que
        romper una red de iones frente a débiles fuerzas de dispersión (\qty{770}{\degreeCelsius} frente a
        \qty{-182}{\degreeCelsius}).
    \end{solu}
\end{exercise}

\begin{exercise}
    \(\star\star\star\) Responde: a) ¿Por qué el \ch{NH3} es muy soluble en agua, pero el \ch{N2} y el \ch{H2} casi no?
    b) Si el enlace B--Cl es polar, ¿por qué la molécula de \ch{BCl3} no lo es? c) ¿Todos los metales son sólidos a temperatura
    ambiente y buenos conductores?
    \begin{solu}
        a) \textcolor{red}{El \ch{NH3} es polar y forma enlaces de hidrógeno con el agua}; el \ch{N2} y el \ch{H2} son apolares.
        b) \textcolor{red}{Porque es triangular plana}: los tres momentos de enlace, a \ang{120}, se anulan. c)
        \textcolor{red}{Todos conducen bien, pero no todos son sólidos}: el mercurio es líquido, y el galio y el cesio funden con
        el calor de la mano (\(\approx \qty{29}{\degreeCelsius}\)).
    \end{solu}
\end{exercise}

\begin{exercise}
    \(\star\star\star\) ¿Conduce la electricidad un hilo de plata? ¿Y el \ch{HBr} disuelto en agua? ¿Y un cristal de \ch{KNO3}?
    ¿Y una disolución de \ch{KNO3}? ¿Y el \ch{KNO3} fundido?
    \begin{solu}
        \textcolor{red}{El hilo de plata, sí} (metal). \textcolor{red}{El \ch{HBr} en agua, sí}: aunque es covalente, en agua se
        disocia en iones \ch{H+} y \ch{Br-}. \textcolor{red}{El cristal de \ch{KNO3}, no}: sus iones no se mueven.
        \textcolor{red}{La disolución y el \ch{KNO3} fundido, sí}: los iones quedan libres.
    \end{solu}
\end{exercise}

\subsection{Sustancias moleculares y fuerzas intermoleculares}

\begin{exercise}
    \(\star\star\star\) Dibuja las estructuras de Lewis del \ch{CBr4}, el \ch{SiH4} y el \ch{SF4}, y determina cuál de ellos no
    tiene geometría tetraédrica.
    \begin{solu}
        \textcolor{red}{El \ch{SF4}}: el S forma cuatro enlaces y le queda un par solitario (octeto expandido, 10 electrones): tiene
        cinco pares alrededor y su forma es de «balancín». El \ch{CBr4} y el \ch{SiH4} tienen cuatro pares enlazantes y son
        tetraédricos.
    \end{solu}
\end{exercise}

\begin{exercise}
    \(\star\star\star\) Dibuja las estructuras de Lewis del \ch{H2S}, el \ch{CO2}, el \ch{Cl2} y el \ch{HCN}, y determina cuál de
    ellos no es lineal.
    \begin{solu}
        \textcolor{red}{El \ch{H2S}}, angular por sus dos pares solitarios. El \ch{CO2} y el \ch{HCN} (sin pares solitarios en el
        C) son lineales, y el \ch{Cl2}, como toda molécula diatómica, también.
    \end{solu}
\end{exercise}

\begin{exercise}
    \(\star\star\star\) ¿Cuál de estas parejas compuesto/geometría es incorrecta? a) \ch{Br2}, lineal. b) \ch{CH2O}, triangular
    plana. c) \ch{O3}, angular. d) \ch{NH4+}, tetraédrica. e) \ch{H2S}, piramidal.
    \begin{solu}
        \textcolor{red}{La e}: el \ch{H2S} tiene dos pares enlazantes y dos solitarios: es angular.
    \end{solu}
\end{exercise}

\begin{exercise}
    \(\star\star\star\) Dibuja las estructuras de Lewis del \ch{BCl3}, el \ch{PCl3} y el \ch{CH2O}, y determina cuál de ellos no
    tiene geometría triangular plana.
    \begin{solu}
        \textcolor{red}{El \ch{PCl3}}: tiene un par solitario y es una pirámide trigonal. El \ch{BCl3} (tres enlaces simples) y el
        \ch{CH2O} (dos simples y un doble, que cuenta como una dirección) son triangulares planos.
    \end{solu}
\end{exercise}

\begin{exercise}
    \(\star\star\star\) A partir del \ch{NH3} y del \ch{NH4+}, determina: a) sus estructuras de Lewis; b) su geometría; c) si la
    molécula de \ch{NH3} es polar; d) por qué las sales de amonio, como el \ch{NH4Cl}, son solubles en agua.
    \begin{solu}
        a) \ch{NH3}: \chemfig{H-\charge{90=\:}{N}(-[6]H)-H}; \ch{NH4+}: el par solitario del N forma un enlace dativo con un
        \ch{H+} (figura~\ref{fig:ud4-dativo}). b) \textcolor{red}{\ch{NH3}, pirámide trigonal; \ch{NH4+}, tetraédrica}. c)
        \textcolor{red}{Sí}, porque es piramidal. d) \textcolor{red}{Son compuestos iónicos}, formados por iones \ch{NH4+} y
        aniones, que el agua hidrata como a los de cualquier sal.
    \end{solu}
\end{exercise}

\begin{exercise}
    \(\star\star\star\) ¿Cuáles de estas parejas de moléculas tienen la misma geometría? Razona la respuesta. a) \ch{GeH4} y
    \ch{CCl4}; b) \ch{AsH3} y \ch{NH3}; c) \ch{Br2} y \ch{F2}; d) \ch{H2S} y \ch{CO2}.
    \begin{solu}
        \textcolor{red}{a), b) y c)}. a) Las dos son tetraédricas (átomo central del grupo 14, cuatro enlaces). b) Las dos son
        pirámides trigonales (grupo 15, tres enlaces y un par solitario). c) Las dos son lineales (diatómicas). d)
        \textcolor{red}{No}: el \ch{H2S} es angular y el \ch{CO2}, lineal.
    \end{solu}
\end{exercise}

\begin{exercise}
    \(\star\star\star\) Teniendo en cuenta sus fuerzas intermoleculares, ordena de menor a mayor punto de ebullición: \ch{Cl2},
    \ch{HF}, \ch{F2} y \ch{HCl}. ¿Te sorprende la posición del \ch{Cl2}?
    \begin{solu}
        \textcolor{red}{\(\ch{F2} < \ch{HCl} < \ch{Cl2} < \ch{HF}\)} (\(-188\), \(-85\), \(-34\) y \qty{20}{\degreeCelsius}). El
        \ch{F2} solo tiene fuerzas de dispersión débiles (molécula pequeña); el \ch{HF}, enlaces de hidrógeno. Puede sorprender que
        el \ch{Cl2}, apolar, hierva por encima del \ch{HCl}, polar: \textcolor{red}{su molécula tiene muchos más electrones
        (34 frente a 18) y sus fuerzas de dispersión superan a las dipolo-dipolo del \ch{HCl}}.
    \end{solu}
\end{exercise}

\begin{exercise}
    \(\star\star\star\) Razona si son correctas estas afirmaciones: a) La energía de las fuerzas de dispersión es menor que la de
    los enlaces de hidrógeno. b) Las fuerzas de dispersión aparecen entre todas las moléculas. c) Al fundir hielo se rompen enlaces
    O--H. d) Los compuestos con enlaces de hidrógeno tienen siempre puntos de ebullición más altos que los que no los tienen.
    \begin{solu}
        a) \textcolor{red}{Correcta en general} para moléculas de tamaño parecido. b) \textcolor{red}{Correcta}: toda molécula
        puede formar dipolos instantáneos. c) \textcolor{red}{Incorrecta}: se rompen parte de los enlaces de hidrógeno entre
        moléculas; los enlaces covalentes O--H siguen intactos. d) \textcolor{red}{Incorrecta}: en moléculas muy grandes las fuerzas
        de dispersión pueden ser más intensas; por ejemplo, el octano (\ch{C8H18}, apolar) hierve a \qty{126}{\degreeCelsius}, más
        que el agua.
    \end{solu}
\end{exercise}
